% Cote 119: offprints, reports, courses, CV and letters (typed and printed part): archive cote 119, site page 191 (PDF page 192).
% typescript. Draft compared: none.
% AI-transcribed by Claude Opus 5.5 (claude-opus-5-5, adaptive thinking, effort high), Message Batches API, batch r119_full (msgbatch_01DwZNm5Y7YsKBWyAN4sFs8w), 2026-11-13; not yet checked by eye.
\scanpage{119/191}
\mspage{- 10 -}

où $\mathrm{Res}(f,z_k)$ désigne le résidu de $f$ en $z_k$ et $I(\gamma,z_k)$ l'indice du che-%
min $\gamma$ par rapport à $z_k$.

\uline{Démonstration} :

D'après la proposition (5.2.) $I(\gamma,z_k) = 0$ pour tout $z_k$ appartenant
à la composante connexe non bornée du complémentaire de l'image de $\gamma$.
Comme les $z_k$ sont isolés, il ne peut y avoir qu'un nombre fini de $z_k$ qui
n'appartiennent pas à cette composante connexe dont le complémentaire est com-%
pact (voir 5.2. fin de la démonstration).

Soit $z_1,z_2,\dots,z_n$ les points $z_k$ pour lesquels $I(\gamma,z_k) \neq 0$.
Considérons la fonction $g$ définie par :
\[ g(z) = f(z) \quad \sum_{k=1}^{n} \frac{\mathrm{Rés}\,(f,z_k)}{z-z_k} \]
Tout revient à démontrer que $\displaystyle\int_\gamma g(z)\,dz = 0$ or la forme différentielle
$g(z)\,dz$ est fermé dans $U$. Elle l'est dans $U$ privé de $z_1,z_2,\dots,z_n$ par le
théorème de Cauchy (4.3) et on voit comme en 3.1. qu'elle admet une primitive
au voisinage de $z_k(1 \leq k \leq n)$. Il ne reste donc plus qu'à appliquer le théo-%
rème 6.8. pour conclure.

\begin{center}\rule{4cm}{0.4pt}\end{center}

\margin{191}
